% Reader question: how does one negative group test avoid individual queries?
% Claim: F=16, P={2} yields nine queries; the negative half of size eight costs
% one query rather than eight singleton queries, under the noiseless oracle.
% Evidence: the exact adaptive-halving procedure in thm:zoom-advantage;
% Q=2*log2(16)+1=9. This is a constructed trace, not a measured workload.
% Grammar: rooted binary query tree, positive path and terminal negatives.
\begin{figure}[H]
\centering
\pdfigurehue{pdteal}%
\begin{tikzpicture}[pd figure,x=1cm,y=1cm]
  \node[pd title,anchor=west] at (0,4.95) {One positive among sixteen candidates};
  \node[pd focus state,minimum width=1.6cm] (R) at (2.0,4.1) {$1$--$16$};
  \node[pd focus state,minimum width=1.6cm] (A) at (2.0,2.95) {$1$--$8$};
  \node[pd artifact,minimum width=1.6cm] (B) at (5.2,2.95) {$9$--$16$};
  \node[pd focus state,minimum width=1.6cm] (C) at (2.0,1.8) {$1$--$4$};
  \node[pd artifact,minimum width=1.6cm] (D) at (5.2,1.8) {$5$--$8$};
  \node[pd focus state,minimum width=1.6cm] (E) at (2.0,.65) {$1$--$2$};
  \node[pd artifact,minimum width=1.6cm] (F) at (5.2,.65) {$3$--$4$};
  \node[pd artifact,minimum width=1.2cm] (G) at (.8,-.6) {$1$};
  \node[pd focus state,minimum width=1.2cm] (H) at (3.2,-.6) {$2$};
  \foreach \a/\b in {R/A,R/B,A/C,A/D,C/E,C/F,E/G,E/H}
    \draw[pd thin arrow] (\a) -- (\b);
  \node[pd label,anchor=west,text width=3.7cm,align=left] at (6.4,2.95)
    {negative: discard eight candidates after one query};
  \node[pd label,anchor=west,text width=3.7cm,align=left] at (6.4,1.25)
    {unfilled groups are negative; coloured groups contain the positive};
  \node[pd title,anchor=west] at (6.4,-.35) {9 queries; flat inspection: 16};
\end{tikzpicture}
\caption{A constructed halving trace with positive item $2$.
Every box is one group query; only positive nonsingletons split.
Negative groups certify all their members at once. The saving counts
noiseless, unit-cost queries, not elapsed time or bytes scanned.}
\label{fig:legible-halving-tree}
\end{figure}
